\documentclass[11pt]{article}
\usepackage[margin=0.85in]{geometry}
\usepackage{amsmath,amssymb,bm}
\setlength{\parindent}{0pt}
\setlength{\parskip}{0.45em}
\newcommand{\Rmat}{\mathcal{R}}
\newcommand{\Vmat}{\mathcal{V}}
\newcommand{\Gmat}{\mathcal{G}}
\newcommand{\Qmat}{\mathcal{Q}}
\newcommand{\Pmat}{\mathcal{P}}

\begin{document}
\begin{center}
{\large \textbf{Manual Derivation: Huggett Discrete-State Problem}}\\[2pt]
ECONS 509 -- Problem Set 1
\end{center}

\textbf{Household problem.}
For current assets $k$ and labor-efficiency state $s\in\{1,2\}$,
\[
v(k,s)=\max_{k'\ge \underline{k}}\left\{u(c)+\beta\sum_{s'=1}^{2}P_{ss'}v(k',s')\right\},
\qquad
c=(1+r^0)k+w^0\epsilon_s\bar l-k',
\]
where $u(c)=c^{1-\sigma}/(1-\sigma)$ and, in the discrete problem,
$k'\in\mathcal K=\{k_1,\ldots,k_N\}$.  The prices are fixed at
\[
w^0=(1-\alpha)z^*\left(\frac{K^0}{L^*}\right)^\alpha,
\qquad
r^0=\alpha z^*\left(\frac{K^0}{L^*}\right)^{\alpha-1}-\delta,
\]
with $L^*=\sum_s\pi_s^*\epsilon_s\bar l$.

\textbf{Discrete-state vector form.}
Let $S=2$ and $M=NS$.  Store the value function as
$\Vmat\in\mathbb R^{N\times S}$, with column $s$ equal to the values on the
asset grid for shock $s$, and stack columns so that
$m(n,s)=(s-1)N+n$.  Thus $\Vmat(:)\in\mathbb R^{M}$.
For current joint state $m(n,s)$ and candidate next-asset index $j$, define
\[
c_{m j}=(1+r^0)k_n+w^0\epsilon_s\bar l-k_j,
\qquad
\Rmat_{m j}=\begin{cases}
\dfrac{c_{m j}^{1-\sigma}}{1-\sigma},&c_{m j}>0,\\[5pt]
-\infty,&c_{m j}\le0.
\end{cases}
\]
Hence $\Rmat\in\mathbb R^{M\times N}$.  Since
$\Pmat\Vmat^T\in\mathbb R^{S\times N}$, the Kronecker product
$(\Pmat\Vmat^T)\otimes\mathbf 1_N$ is $M\times N$.  The Bellman equation is
\[
\boxed{\quad
\Vmat(:)=\max_{j=1,\ldots,N}
\left[\Rmat+\beta(\Pmat\Vmat^T)\otimes\mathbf 1_N\right],
\quad}
\]
where the maximum is row by row.  The policy
$\Gmat\in\{1,\ldots,N\}^{N\times S}$ stores the maximizing asset index,
\[
\Gmat_{ns}=\arg\max_j\left\{\Rmat_{m(n,s),j}
+\beta\sum_{s'}P_{ss'}\Vmat_{j s'}\right\}.
\]
Given $\Gmat$, define the policy transition matrix
$\Qmat\in\mathbb R^{M\times M}$ by
\[
\Qmat_{m(n,s),m(j,s')}=
\begin{cases}
P_{ss'},&j=\Gmat_{ns},\\
0,&\text{otherwise}.
\end{cases}
\]
If $\Rmat^*\in\mathbb R^M$ collects flow utility at the chosen policy, fixed-policy
evaluation satisfies
\[
\Vmat(:)=\Rmat^*+\beta\Qmat\Vmat(:),
\qquad
(I-\beta\Qmat)\Vmat(:)=\Rmat^*.
\]

\textbf{Continuous-choice Euler equation and borrowing constraint.}
For the continuous-choice version, write the constraint as
$k'-\underline{k}\ge0$ with multiplier $\lambda\ge0$.  Using the resource
constraint to substitute for consumption, the Lagrangian is
\[
\mathcal L=u\!\left((1+r^0)k+w^0\epsilon_s\bar l-k'\right)
+\beta\sum_{s'}P_{ss'}v(k',s')+\lambda(k'-\underline{k}).
\]
The first-order condition for $k'$ and the envelope condition are
\[
-u'(c)+\beta\sum_{s'}P_{ss'}v_k(k',s')+\lambda=0,
\qquad
v_k(k,s)=(1+r^0)u'(c).
\]
Combining them gives the KKT Euler condition
\[
\boxed{\quad
u'(c)=\beta(1+r^0)\sum_{s'}P_{ss'}u'(c'(s'))+\lambda,
\quad}
\]
with
\[
\lambda\ge0,\qquad k'-\underline{k}\ge0,\qquad
\lambda(k'-\underline{k})=0.
\]
For CRRA utility, $u'(c)=c^{-\sigma}$.  Therefore, whenever the borrowing
constraint is slack ($k'>\underline{k}$), $\lambda=0$ and
\[
c^{-\sigma}=\beta(1+r^0)\sum_{s'}P_{ss'}\,[c'(s')]^{-\sigma}.
\]

\textbf{Euler residual used in the accuracy check.}
Let the discrete policy be $g(k,s)$ and define policy consumption
\[
c_g(k,s)=(1+r^0)k+w^0\epsilon_s\bar l-g(k,s).
\]
At a slack state, put $k'=g(k,s)$ and, for each $s'$, let
$k''=g(k',s')$ so that $c_g(k',s')$ is next-period consumption.  I use the
consumption-equivalent relative Euler residual
\[
\boxed{\quad
e(k,s)=1-
\frac{\left[\beta(1+r^0)\sum_{s'}P_{ss'}c_g(k',s')^{-\sigma}\right]^{-1/\sigma}}
{c_g(k,s)}.
\quad}
\]
The continuous-choice Euler equation implies $e(k,s)=0$.  In Question 2(h) I
report $|e(k,s)|$ only where $g(k,s)>\underline{k}$; binding states obey the
KKT inequality rather than the Euler equality.

\end{document}
